\section{An arithmetic arrangement and a distinguished surface}\label{sec:geometry}

We construct a smooth proper effective orbifold $\cS$ with projective
coarse space and a morphism $f:Y\to\cS$, where $Y$ is obtained from a
simple abelian surface by point blowups.  The image of $f$ will lie in
the ordinary locus of $\cS$.  The arrangement constructed here provides
both the relations and the local covering spaces used to prove that
$f_*\pi_1(Y)$ is infinite cyclic.  The fundamental-group calculation is
made in Theorem~\ref{thm:orbifold-cyclic}.

\subsection{Two pencils and their coupling hyperplanes}\label{geo:prototype}

The two pencils will use the same six marked directions on $\PP^1$.
We first choose these markings so that their double cover has simple
Jacobian, and then choose the arithmetic ball containing the pencils.
Choose six distinct real algebraic numbers $\lambda_1,\ldots,\lambda_6$ with
\begin{equation}\label{geo:slopes}
 |\lambda_i|<0.01\quad(1\le i\le4),
 \qquad \lambda_5<-4,\quad \lambda_6>4,
\end{equation}
and such that the smooth double cover
\begin{equation}\label{geo:curve}
 \pi:C\longrightarrow\PP^1,
 \qquad y^2=\prod_{i=1}^6(t-\lambda_i),
\end{equation}
has simple Jacobian.  We call the first four slopes \emph{red} and the
last two \emph{blue}.  These requirements can be imposed simultaneously,
as follows.

Start with six disjoint real open intervals satisfying
\eqref{geo:slopes}.  Monic real polynomials with one root in each interval
form a nonempty open set in coefficient space.  At three sufficiently
large distinct primes prescribe squarefree degree-six reductions with
factor degrees
\[
 (6),\qquad(5,1),\qquad(2,1,1,1,1).
\]
Weak approximation produces a monic polynomial in $\Q[t]$, integral at
these primes, with these reductions and with its real coefficients in
the chosen open set.  Its Galois group contains a six-cycle, a five-cycle,
and a transposition, by the Frobenius cycle criterion.  The six-cycle
makes the group transitive.  The stabilizer of the fixed point of the
five-cycle is transitive on the remaining five points, so the group is
two-transitive.  Conjugating its transposition then gives every
transposition, and the Galois group is $S_6$.  Let $F$ be its splitting
field.  Then $F$ is totally real, $F\ne\Q$, and all six slopes belong
to $F$.  Zarhin's theorem gives
$\End(\Jac(C))=\Z$ over an algebraic closure, hence $\Jac(C)$ is simple
\cite[Theorem~2.1]{Zarhin}.

Use the inclusion $F\subset\R$ containing the chosen roots as the
distinguished real embedding.  Choose a CM extension $E/F$ and an
element $a\in F$ which is positive at this embedding and negative at
every other real embedding; weak approximation supplies $a$.  On $E^5$ put
\begin{equation}\label{geo:hermitian}
 h=\operatorname{diag}(1,1,1,1,-a).
\end{equation}
At the distinguished embedding the negative lines form complex
hyperbolic space of dimension four.  In the affine chart with last
homogeneous coordinate one, write this ball as
\[
 \BB=\{(z,w)\in\C^2\times\C^2:|z|^2+|w|^2<a\},
 \qquad o=(0,0).
\]
The metric is normalized so that in the unit-radius ball a point at
distance $r$ from the origin has Euclidean radius $\tanh r$.

Choose $c\in F$ with
\begin{equation}\label{geo:coupling-size}
                 0.85<c/\sqrt a<0.9.
\end{equation}
The \emph{$z$-family} consists of the six slope hyperplanes and two
coupling hyperplanes
\[
 H^z_i=\{z_2=\lambda_i z_1\}\cap\BB\quad(1\le i\le6),
 \qquad K^z_\pm=\{z_1=\pm c\}\cap\BB.
\]
Define the $w$-family by the same equations in $w$.  Let $\cD$ denote
these sixteen labeled complex hyperbolic hyperplanes.  They are all
defined over $E$ and are preserved, with their labels permuted, by
\begin{equation}\label{geo:sigma}
                    \sigma(z,w)=(w,-z).
\end{equation}
This is an order-four isometry represented by an integral unitary
matrix.  The two pencils have centers
\[
 L_z=\{z=0\}\cap\BB,
 \qquad L_w=\{w=0\}\cap\BB.
\]

The intersection pattern is elementary but important.  The $z$ slopes
all meet along $L_z$.  On $K^z_\pm$, a slope of value $\lambda$ meets the
ball exactly when $c^2(1+|\lambda|^2)<a$.  Thus each coupling hyperplane
meets the four red slopes, in distinct subspaces disjoint from $L_z$,
and meets neither blue slope.  The two $z$ coupling hyperplanes are
disjoint.  These statements hold also in the $w$-family.  A $z$ coupling
hyperplane and a $w$ coupling hyperplane are disjoint because $2c^2>a$.
Every intersection involving the two families has a local product
description in the $z$ and $w$ variables.  In particular, the phrase
``product'' here concerns local analytic coordinates, rather than a
product decomposition of the hyperbolic metric.

Figure~\ref{fig:arrangement} shows the incidences in one family; the
second family is obtained by using the other coordinate pair.
\begin{figure}[htbp]
\centering
\begingroup
\definecolor{arrred}{RGB}{165,58,37}
\definecolor{arrblue}{RGB}{31,87,145}
\begin{tikzpicture}[x=1cm,y=1cm,>=Latex,
  every node/.style={font=\small},
  redline/.style={arrred,line width=.8pt},
  blueline/.style={arrblue,dashed,line width=.8pt},
  redpoint/.style={circle,fill=arrred,draw=arrred,inner sep=1.8pt},
  bluepoint/.style={rectangle,fill=arrblue,draw=arrblue,inner sep=1.8pt}]
  \node at (2.55,4.55) {The six marked directions};
  \draw[->,black!65] (0,2.6)--(5.1,2.6);
  \node[right] at (5.1,2.6) {$\lambda$};
  \draw[black!45,rounded corners=9pt] (1.35,2.12) rectangle (3.75,3.08);
  \foreach \x in {1.7,2.25,2.8,3.35}
    \node[redpoint] at (\x,2.6) {};
  \node[bluepoint] at (.35,2.6) {};
  \node[bluepoint] at (4.75,2.6) {};
  \node[align=center,arrred] at (2.55,3.6)
    {four red points\\$|\lambda_i|<0.01$};
  \node[align=center,arrblue] at (2.55,1.45)
    {two blue points: $|\lambda_i|>4$};
  \node[align=center] at (2.55,.45)
    {$C\longrightarrow\mathbb P^1_{\lambda}$\\double cover branched at these points};

  \begin{scope}[shift={(9.15,2.25)}]
    \node at (0,2.3) {One pencil and its couplings};
    \draw[black!40,line width=.6pt] (0,0) circle (2.05);
    \foreach \angle in {-13,-4,4,13}
      \draw[redline] (\angle:2.05)--({\angle+180}:2.05);
    \foreach \angle in {78,102}
      \draw[blueline] (\angle:2.05)--({\angle+180}:2.05);
    \draw[line width=1pt] (-1.78,-1.017)--(-1.78,1.017);
    \draw[line width=1pt] (1.78,-1.017)--(1.78,1.017);
    \foreach \angle in {-13,-4,4,13}{
      \fill[arrred] (1.78,{1.78*tan(\angle)}) circle (1.9pt);
      \fill[arrred] (-1.78,{-1.78*tan(\angle)}) circle (1.9pt);
    }
    \fill (0,0) circle (2pt);
    \node[fill=white,inner sep=1.5pt,below=8pt] at (0,0)
      {$L_z$};
    \node[align=center,below=4pt,fill=white,inner sep=1.5pt] at (-1.78,-1.017)
      {$z_1=-c$};
    \node[align=center,below=4pt,fill=white,inner sep=1.5pt] at (1.78,-1.017)
      {$z_1=c$};
    \node[arrblue,fill=white,inner sep=1pt] at (.86,1.64) {blue};
    \node[arrred,fill=white,inner sep=1pt] at (.92,.57) {red};
    \node at (0,-2.55) {$z_2=\lambda_i z_1$};
  \end{scope}
\end{tikzpicture}
\endgroup
\caption{The finite arrangement in one coordinate pair. Left: a
  schematic placement of the six branch points on the real direction
  axis, grouped into four red circles and two blue squares. Right:
  the slice $w=0$, $z_1,z_2\in\mathbb R$ of the $z$ pencil in the
  ball; solid red slopes meet both coupling
  hyperplanes, whereas dashed blue slopes meet neither. The red
  angular separation is exaggerated to show these incidences.
  The center represents the slice of $L_z$. The $w$ family is an
  identical construction in the other coordinate pair; the full
  arrangement lies in complex dimension four.}
\label{fig:arrangement}
\end{figure}


\subsection{Exact models in congruence quotients}\label{geo:separation}

The arithmetic lattice theorem of Borel--Harish-Chandra
\cite[Corollary~12.4]{BorelHarishChandra1962}, applied to
\eqref{geo:hermitian}, gives a
cocompact arithmetic lattice of simplest type in $\mathrm{PU}(4,1)$.
Indeed all the other archimedean factors are compact, and the form is
anisotropic over $E$: a nonzero $E$-rational isotropic vector would
remain isotropic at a definite conjugate embedding.  Torsion-free
principal congruence subgroups therefore give compact complex
hyperbolic manifolds, which are projective by the Baily--Borel theorem
\cite[Theorem~10.11]{BailyBorel1966}.  One can
also obtain projectivity from the positive canonical bundle of a
compact ball quotient.  We use principal congruence subgroups normalized
by $\sigma$, and may take their levels arbitrarily deep.  These standard
arithmetic facts are recalled in \cite[Section~3.1]{StoverToledo}.

Fix positive normal vectors $n_i\in E^5$ for the members of $\cD$,
scaled to have integral coordinates.  We always label a translate of
the $i$th hyperplane by $i$, and use the prescribed normal $\gamma n_i$
for its translate by $\gamma$.  Congruence subgroups may be chosen to
avoid the finite scalar kernel of the projective action.

\begin{proposition}[Exact local arrangements]\label{prop:arrangement-local}
Given $W>0$ and $R_0>0$, there is a torsion-free principal congruence
subgroup $\Gamma_0$, normalized by $\sigma$, with the following
properties.  Let $\mathcal A_0$ be the union of all $\Gamma_0$-translates
of the labeled hyperplanes of $\cD$.
\begin{enumerate}
\item The arrangement $\mathcal A_0$ is locally finite.  The hyperplanes
meeting any ball of radius $W$ have distinct labels and are carried to
the corresponding subarrangement of $\cD$ by a holomorphic isometry
induced by a unitary map sending each prescribed translated normal to
its prototype normal.
\item The injectivity radius of $\Gamma_0\backslash\BB$ is greater than
$R_0$.
\item In the ball of radius $R_0$ about $o$, the arrangement is exactly
$\cD$ restricted to that ball.  The analogous assertion holds in every
$\Gamma_0$-translate of this ball.
\end{enumerate}
All these conclusions concerning the lifted arrangement remain valid
on any subsequent finite cover of $\Gamma_0\backslash\BB$, provided one
uses the full inverse image of the arrangement.
\end{proposition}

\begin{proof}
First fix a radius $W$ and a finite collection of hyperplanes meeting a
$W$-ball.  At its viewpoint, represented by a unit negative
vector $x$, the distance to the hyperplane with positive normal $n$
satisfies
\[
  \sinh d(x,H_n)=\frac{|h(n,x)|}{\sqrt{h(n,n)}}.
\]
For a normal of one of our finitely many fixed lengths, intersection
with the $W$-ball bounds $|h(n,x)|$.  Write $n=n_++\alpha x$ with
$n_+\in x^\perp$.  Since
\[
 h(n_+,n_+)=h(n,n)+|\alpha|^2,
\]
the positive component is bounded as well.  Cauchy--Schwarz in
$x^\perp$ now bounds every Gram entry of two normals visible from this
ball, by a constant depending only on $W$ and their labels.

At every other archimedean place, definiteness and the fixed normal
lengths bound the same Gram entries.  All entries belong to one fixed
fractional ideal of $E$.  Its image under the archimedean embeddings
is a lattice, so only finitely many Gram entries satisfy these bounds.
If $\gamma,\gamma'$ are in a sufficiently deep integer principal
congruence subgroup, then
\[
 h(\gamma n_i,\gamma'n_j)\equiv h(n_i,n_j)
\]
modulo the level, in this fractional ideal.  Choose the level to
separate all the finitely many possible nonzero differences.  We obtain
exact equality:
\begin{equation}\label{geo:gram}
            h(\gamma n_i,\gamma'n_j)=h(n_i,n_j)
\end{equation}
whenever the corresponding hyperplanes are visible together.

Here equality of Gram matrices gives the asserted isometry even for
dependent collections.  Every $E$-rational subspace of $E^5$ is
nondegenerate for $h$: a nonzero radical, being defined by linear
equations over $E$, would persist at a definite embedding.  Consequently
the kernel of a Gram matrix is exactly the relation space of its
vectors.  Equality \eqref{geo:gram} thus defines an isometry between
their spans, which extends to a unitary isometry of the ambient
Hermitian space.  Its projectivization is a holomorphic ball isometry.
Two visible copies of the same label must in fact have the same normal:
their difference has norm zero by \eqref{geo:gram}, and anisotropy over
$E$ makes that difference zero.  Two different labels cannot coincide,
since this would contradict equality of their relation spaces with
those of the distinct prototypes.

Local finiteness follows from the same boundedness argument applied to
normals themselves at a fixed viewpoint.  The normals have integral
coordinates; their distinguished components are bounded by the distance
bound and their conjugate components by definiteness.  Only finitely
many can meet a fixed bounded ball.

To obtain exact agreement at $o$, choose an auxiliary ball about $o$
large enough to contain the prescribed $R_0$-ball and to meet every
prototype hyperplane.  Apply the Gram argument to this ball.  The
prototype normals span $E^5$: the slope normals span the four positive
coordinate directions, and a coupling normal supplies the last
homogeneous direction.  An additional translated normal visible in
this ball has, by \eqref{geo:gram}, the same inner products with this
spanning set as its prototype.  It therefore equals that prototype.
This proves the third assertion.

Finally fix a torsion-free congruence subgroup in advance and a compact
fundamental set for its action on $\BB$.  A group element moving a point
of this set a bounded distance belongs to a finite set.  A nested
sequence of normal principal congruence subgroups has trivial
intersection, after the scalar kernel has been removed.  Deep enough
levels avoid every nonidentity element of the finite set.  Normality
allows an arbitrary viewpoint to be moved into the compact fundamental
set without changing subgroup membership.  This proves the required
uniform lower bound on injectivity radius.  Increasing the level
accommodates all the preceding requirements simultaneously.  Passing
to a subgroup preserves them when the arrangement is pulled back in
full.
\end{proof}

We next record the consequences for the divisors and centers that we
will blow up.  In $M_0=\Gamma_0\backslash\BB$, each label gives a closed
embedded smooth totally geodesic divisor, possibly disconnected.  Its
lifts are locally finite and have disjoint sheets by
Proposition~\ref{prop:arrangement-local}.  The divisors are algebraic
because $M_0$ is projective.

Within one family the intersection centers have complex codimension
two.  A center is either a common intersection of at least two slopes,
or an intersection of a red slope with a coupling hyperplane.  At a
slope center only a subset of the six slopes may occur.  The subset is
constant along a connected center: if another slope meets it, the
local model shows that it contains the intersection, and the
containment extends along its lifts.  Distinct centers within one
family are disjoint; otherwise all participating hyperplanes would be
visible near a point of intersection, contradicting the prototype
pattern.  Coincident centers are counted once.  Thus these centers are
closed embedded smooth subvarieties.  Centers from different families
have the product structure already described.  These assertions apply
to the full pulled-back arrangement after any finite cover.

\subsection{Two global sign covers}\label{geo:signs}

For the later local models we need global parity characters on the
arrangement complement.  We obtain these from the following theorem
of Stover--Toledo, in the form stated by Llosa Isenrich--Py
\cite[Theorem~4.2]{LlosaPy}; see also \cite{StoverToledo}.

\begin{theorem}[Virtual cyclic ramified covers]\label{geo:stover-toledo}
Let $M_0$ be a compact ball quotient of complex dimension at least two
by a torsion-free congruence arithmetic lattice of simplest type.  Let
$D$ be a nonempty totally geodesic divisor whose components are smooth,
embedded, and pairwise disjoint.  For every integer $d\ge2$ there is a
finite unramified cover $M_1\to M_0$ such that $M_1$ admits a cyclic
cover of degree $d$, branched over the full inverse image of $D$.
\end{theorem}

Apply Theorem~\ref{geo:stover-toledo} with $d=2$ to each of the sixteen
label divisors separately, on the initial congruence quotient $M_0$.
Each application satisfies the disjoint-component hypothesis.  We do
not apply the theorem to the union of the intersecting divisors.  Take
a common finite unramified cover of the resulting base covers.  The
sum, in $\Z/2$, of the eight characters for the $z$ labels gives a
character $\chi_z$ on the complement of the $z$ arrangement.  It has
value one on a meridian of every $z$ hyperplane.  Pull it back to the
complement of both families.  The same procedure gives a $w$ character.
Subsequent finite covers need not be congruence: every use of
Theorem~\ref{geo:stover-toledo} has already taken place on $M_0$.

We arrange equivariance before fixing the signs.  First intersect the
finitely many $\sigma$-conjugates of the subgroup defining our common
cover, so that $\sigma$ acts on it.  Since $\sigma^2$ preserves the
$z$-family, the character
\[
                 \delta=\chi_z+\sigma^{2*}\chi_z
\]
is zero on all meridians.  The kernel of the map from the fundamental
group of a divisor complement to that of the ambient smooth variety
is normally generated by meridians.  Hence $\delta$ factors through
the fundamental group of the compact ball quotient.  Pass to a further
finite cover by intersecting the kernel of this character with its
$\sigma$-conjugates.  On that cover the pulled-back $z$ character
satisfies $\sigma^{2*}\chi_z=\chi_z$.  Define
\begin{equation}\label{geo:sign-equivariance}
                  \chi_w=\sigma^*\chi_z.
\end{equation}
Then $\sigma$ exchanges the two characters.  They have the prescribed
meridian values in their respective families, and each is defined
using only that family.

Write $M=\Gamma\backslash\BB$ for this final finite cover, still with
the full pulled-back arrangement.  The image of $o$ in $M$, also denoted
$o$, is fixed by $\sigma$.  It is an isolated component even of the
fixed locus of $\sigma^2$, because the derivative of $\sigma^2$ at $o$
is $-\id$.  No assertion of this kind is needed for the other fixed
loci of the symmetry on $M$.

\subsection{Filling the arrangement}\label{geo:fill}

Blow up all intersection centers within the $z$-family, and then the
transforms of those within the $w$-family.  Product coordinates show
that the operations for the two families commute.  The branch rule
for the $z$ character is as follows.  Every strict transform of a
$z$ hyperplane is a branch divisor of order two.  If $k$ such
hyperplanes pass through a center, its exceptional meridian is the
product of their meridians, and therefore has sign $k\bmod2$.
Accordingly the exceptional divisor is branched precisely when $k$
is odd.  In that case blow up each meeting of this exceptional divisor
with a strict transform of a hyperplane.  The two meeting divisors
both have sign one; their new exceptional divisor has sign zero.
The branch components within this family are now smooth and disjoint.
Apply the identical rule to the $w$ character.  This prescription
includes partial pencils.  In particular an ordinary red--coupling
pair has $k=2$, so its first exceptional divisor is unbranched.

Let $\widehat M$ denote the resulting smooth projective base.  Extend
the two sign covers from the complement by normalization over
$\widehat M$, and take them together.  Finite topological covers of
smooth complex algebraic varieties algebraize, and normalization in
the resulting finite extensions gives finite algebraic covers; thus
this procedure is projective.  Locally along a branch divisor the
normalization is the smooth double-cover model $t=u^2$.  Branch
divisors within a family are disjoint, and across the families the
local equations and the two signs are independent.  Their combined
normalization is therefore smooth, including at crossings, where it
has product equations $(t_1,t_2)=(u_1^2,u_2^2)$.  We obtain a smooth
projective variety
\begin{equation}\label{geo:V}
                 V\longrightarrow\widehat M
\end{equation}
with deck group $(\Z/2)^2$.  It is connected: a meridian at a general
point of a $z$ divisor has sign $(1,0)$, and a meridian at a general
point of a $w$ divisor has sign $(0,1)$, so the combined monodromy is
surjective.

Near $o$ all six slopes in each pencil occur and the coupling
hyperplanes are absent after the neighborhood has been made small.
The blowups are the product of the blowups of the two coordinate
planes at their origins.  The exceptional fiber of $\widehat M\to M$
is $\PP^1\times\PP^1$.  Each exceptional line has six branch markings
and has even meridian sign.  Its inverse image in $V$ is consequently
\begin{equation}\label{geo:Fstar}
                            F_*=C\times C.
\end{equation}
Put
\[
       L=\pi^*\mathcal O_{\PP^1}(-1),
       \qquad N=N_{F_*/V}
          =\operatorname{pr}_1^*L\oplus\operatorname{pr}_2^*L.
\]
In fact a neighborhood of $F_*$ has the product line-bundle description
provided by these two tautological lines.  Indeed, in each factor, after
the unbranched radial disks have been filled, their disk bundle retracts onto its
zero section.  Away from the six marked directions its cover is thus
pulled back from the direction line.  Normalization gives the
six-branch cover in that direction variable, with the radial line
bundle pulled back unchanged.

The equivariance \eqref{geo:sign-equivariance} makes the kernel of the
combined sign character invariant, so $\sigma$ lifts on the complement.
The equivariant blowups and normalization extend the lift to an
automorphism $s$ of $V$.  Its restriction to $F_*$ exchanges the two
copies of $C$, possibly followed by either hyperelliptic deck
involution.  Composing with a deck transformation makes this
restriction the pure swap.  To determine the lift in the whole local
model, note that $u,v$ are actual tautological vectors: the covering
map in these coordinates is
\[
 (p,q;u,v)\longmapsto(\pi(p),\pi(q);u,v).
\]
The pure swap determines the lifted direction maps on the connected
local cover.  Since the base map sends the radial vectors to $(v,-u)$,
there is no further radial multiplier.  Thus
\begin{equation}\label{geo:line-action}
             s(p,q;u,v)=(q,p;v,-u).
\end{equation}
Its fourth power is a deck transformation fixing $F_*$ pointwise,
so $s^4=1$; a nonidentity element of $(\Z/2)^2$ does not fix $C\times C$
pointwise.  Its square is the identity on $F_*$ and is multiplication
by $-1$ on both normal lines.

\subsection{An ordinary neighborhood of the distinguished fiber}
\label{geo:resolution}

The quotient stack $[V/\langle s\rangle]$ is a smooth proper orbifold
with projective coarse space, but a surface mapping into its distinguished
fiber would
meet stabilizers.  We now replace just that fiber by an ordinary
smooth resolution.  Retaining the quotient stack elsewhere allows
all other fixed loci to be left in place.

First blow up $F_*$ in $V$, and take the coarse quotient by
$\langle s^2\rangle$.  The exceptional divisor is
\begin{equation}\label{geo:P}
                          P=\PP_{F_*}(N),
\end{equation}
where projectivization parametrizes lines.  Before taking the quotient,
$s^2$ acts trivially on $P$ and as $t\mapsto-t$ on its transverse
coordinate.  The quotient is thus smooth near $P$, with transverse
coordinate $t^2$.

The residual involution induced by $s$ has fixed points near $P$ only
in $P$: away from the fiber, a fixed point would project to a
nontrivial fixed point of the symmetry near $o$ in $M$.  Its fixed
points in $P$ project to the diagonal of $C\times C$.  On the
projective normal line over that diagonal, it acts by
\[
                       [u:v]\longmapsto[v:-u].
\]
There are exactly two eigensections, with eigenvalues $i$ and $-i$
before projectivization.  These give two disjoint smooth fixed curves.
On the three normal directions to either curve the involution acts
by $-1$: one direction is antisymmetric in the base, one is tangent
to the projective fiber, and the transverse coordinate to $P$ has
eigenvalue $(\pm i)^2=-1$.  Blow up both fixed curves and take the
coarse quotient by the residual involution.  This quotient is smooth
near the distinguished fiber.  Indeed, for the linear local model
\[
              (x,t_1,t_2,t_3)\longmapsto(x,-t_1,-t_2,-t_3),
\]
blowing up the fixed curve $\{t_1=t_2=t_3=0\}$ makes the action a
reflection in the coordinate normal to the exceptional divisor, and
the coarse quotient is smooth.  Finite-order holomorphic actions are
locally linearizable, so this model applies.

All centers just used are closed and lie over $o$.  The blowups and
finite coarse quotients can therefore be performed globally,
producing a projective coarse space $S_0$ which modifies
$V/\langle s\rangle$ only over $o$.  The intermediate coarse spaces
may retain quotient singularities elsewhere; the preceding calculation
establishes smoothness on a neighborhood of the distinguished fiber.
To form $\cS$, use this smooth scheme near that fiber and use
$[V/\langle s\rangle]$ away from it.  This is an algebraic gluing.
More explicitly, in $M$ remove all components of the nontrivial fixed
loci other than the isolated point $o$, and take the resulting invariant
Zariski neighborhood.  Outside $o$ in this neighborhood the action is
free, so both constructions restrict to the same ordinary quotient.
Their gluing is a smooth connected effective orbifold $\cS$ with
coarse space $S_0$.  The map $\cS\to S_0$ is separated and proper:
these properties are local on $S_0$, where it is either the coarse map of a finite-group
quotient stack or the identity.  Since $S_0$ is projective,
$\cS$ is separated and proper over $\C$.  In particular the whole distinguished
fiber lies in the ordinary smooth locus of $\cS$.

\subsection{The surface mapping into the ordinary fiber}\label{geo:surface}

Let $A=\Jac(C)$.  For a genus-two curve, the Abel map
\[
        \Sym^2C\longrightarrow\Pic^2(C)\simeq A
\]
is the blowup at the canonical class, after choosing the displayed
translation.  To recall the geometry, Riemann--Roch gives a unique
effective divisor in every degree-two class except $K_C$; the latter
has the pencil $E_C=|K_C|\simeq\PP^1$.  For the degree-two quotient
$q:C\times C\to\Sym^2 C$, the inverse image of $E_C$ is the graph
$\Gamma_\iota$ of the hyperelliptic involution.  The map $q$ is generically
unramified along this graph, so $q^*E_C=\Gamma_\iota$ as divisors;
ramification at its six diagonal points does not change the generic
multiplicity.  The projection formula gives
\[
  2E_C^2=(q^*E_C)^2=\Gamma_\iota^2=2-2g(C)=-2.
\]
The Abel map is therefore a birational
morphism between smooth surfaces with one exceptional $(-1)$-curve,
and is the stated point blowup.

Choose a nonzero rational section $\tau$ of $L$.  On the open subset
of $C\times C$ where both values are defined and nonzero, put
\begin{equation}\label{geo:rational-section}
                 (p,q)\longmapsto[\tau(p),i\tau(q)]\in P.
\end{equation}
This section is equivariant for swapping the two factors and the
involution on $P$.  Indeed \eqref{geo:line-action} gives
\[
       [\tau(p),i\tau(q)]\longmapsto[i\tau(q),-\tau(p)]
                            =[\tau(q),i\tau(p)].
\]
Away from the diagonal it avoids the two fixed curves, and thus gives
a rational map from $\Sym^2C$ to the smooth resolution constructed
above.  All the bundles and maps in this description are algebraic:
the normal-bundle identification is algebraic as well as analytic,
by GAGA on the projective fiber.  Resolve the rational map by a finite
sequence of point blowups of its smooth projective source, using the
projective coarse space as target.  The resolved image remains in its
closed distinguished fiber, where the coarse space agrees with the
ordinary locus of $\cS$.  We therefore obtain a morphism
\begin{equation}\label{geo:Ymap}
                        f:Y\longrightarrow\cS
\end{equation}
whose image lies entirely over $o$.  Thus $Y$ is an iterated point
blowup of the simple abelian surface $A$, and $f(Y)$ lies in the
ordinary locus.

\begin{proposition}[The geometric construction]\label{prop:orbifold-surface}
After any prescribed finite lower bounds on the arrangement radii and
injectivity radius, the construction above gives a smooth connected
proper effective orbifold $\cS$ with projective coarse space, a simple
abelian surface $A=\Jac(C)$, an iterated point blowup $Y\to A$, and a
morphism $f:Y\to\cS$ with image in its ordinary locus.  The distinguished
fiber is obtained from the product $C\times C$, its normal bundle
$\operatorname{pr}_1^*L\oplus\operatorname{pr}_2^*L$, and the action
\eqref{geo:line-action} by the explicit blowups and coarse quotients
above.  Away from that fiber the orbifold is $[V/\langle s\rangle]$,
with $V$ given by the two parity covers and the filling rules of
Section~\ref{geo:fill}.
\end{proposition}

All parameters used in the finite arrangement have now been fixed.
The depth of $\Gamma_0$ has not: Proposition~\ref{prop:arrangement-local}
permits its choice after the constants required for the finite-model
graphs in the next section.  Passing to the finite covers which produce
the signs preserves these estimates.  This order of choices will
allow the local calculations to establish
\[
                    \im(f_*:\pi_1(Y)\to\pi_1(\cS))\simeq\Z.
\]
